\documentclass[12pt]{article}
\usepackage{amsmath,amssymb,geometry,enumitem}
\geometry{margin=1in}
\setlength{\parskip}{0.6em}
\setlength{\parindent}{0pt}

\begin{document}


% ============================================================
\section*{Problem 1 (18 Marks)}
% ============================================================

Johny has \(T=100\) hours per week to allocate between leisure \(L\) and work \(h\), where \(h=T-L\). His wage is \(w=\$20\) per hour and non-labour income is \(\$200\).

\begin{enumerate}[label=(\alph*)]

\item \textbf{(4 marks)}  
Write Johny's weekly budget constraint in terms of leisure \(L\) and income \(C\).  
Identify:
\begin{itemize}
    \item the intercept on the income axis,
    \item the intercept on the leisure axis,
    \item the slope of the budget line.
\end{itemize}

\item \textbf{(4 marks)}  
The government introduces a guaranteed income program:
\[
G = 600,\quad t = 0.5
\]

\begin{itemize}
    \item Write Johny’s new disposable income equation while benefits are positive.
    \item Compute the effective net wage.
    \item Find the break-even level of weekly earnings.
    \item Compute the corresponding hours worked at the break-even point.
\end{itemize}

\item \textbf{(3 marks)}  
Johny initially works 40 hours per week. After the program is introduced, he works 30 hours.

\begin{itemize}
    \item Calculate the change in labour supply.
    \item Explain whether this change is consistent with income and substitution effects.
\end{itemize}

\item \textbf{(3 marks)}  
The working-age population is 5,000,000. The labour force participation rate is 70\%, and the unemployment rate is 6\%.

\begin{itemize}
    \item Compute the labour force.
    \item Compute the number of unemployed workers.
    \item Compute the number of employed workers.
\end{itemize}

\item \textbf{(4 marks)}  
Suppose:
\begin{itemize}
    \item 60,000 unemployed workers find jobs,
    \item 40,000 employed workers lose jobs and begin searching,
    \item 30,000 workers leave the labour force.
\end{itemize}

Compute:
\begin{itemize}
    \item the new number of unemployed workers,
    \item the new labour force,
    \item the new unemployment rate.
\end{itemize}

\end{enumerate}

% ============================================================
\section*{Problem 2 (17 Marks)}
% ============================================================

A firm faces the following labour market conditions:

\begin{itemize}
    \item Labour demand schedule:  
    \[
    W = 50 - 0.5N
    \]
    \item Labour supply schedule:  
    \[
    W = 10 + 0.5N
    \]
\end{itemize}

\begin{enumerate}[label=(\alph*)]

\item \textbf{(4 marks)}  
Find the \textbf{competitive equilibrium}:
\begin{itemize}
    \item equilibrium wage,
    \item equilibrium employment.
\end{itemize}

\item \textbf{(4 marks)}  
Suppose the firm is a monopsonist and hires 20 workers.

\begin{itemize}
    \item Compute the wage paid to workers.
    \item Compare this wage to the competitive wage.
    \item Explain why monopsony leads to lower wages.
\end{itemize}

\item \textbf{(3 marks)}  
A minimum wage of \(\$30\) is introduced.

\begin{itemize}
    \item Compute the number of workers willing to work at this wage.
    \item Compute the number of workers demanded at this wage.
    \item Determine actual employment.
\end{itemize}

\item \textbf{(3 marks)}  
A country admits 400,000 immigrants:
\begin{itemize}
    \item Economic: 200,000
    \item Family: 120,000
    \item Refugees: 80,000
\end{itemize}

\begin{itemize}
    \item Compute the percentage share of each category.
    \item If economic immigrants increase by 10\%, compute the new total number of immigrants.
\end{itemize}

\item \textbf{(3 marks)}  
Johny is considering education that costs \(\$12{,}000\) today and increases earnings by \(\$3{,}000\) per year for 8 years. The discount rate is 5\%.

Use:
\[
PV = A\left(\frac{1-(1+r)^{-T}}{r}\right)
\]

\begin{itemize}
    \item Compute the present value of benefits.
    \item Compute the net present value (NPV).
    \item Should Johny invest? Briefly explain.
\end{itemize}

\end{enumerate}



\end{document}